% Unnumbered post29 research note. No promoted statement or counting claim.
\subsection*{Canonical order saturation and the exact reduction-cycle gap}

Use the original comparable-cycle model and its prescribed pullback to
the subdivided root model. Let $A$ be a fixed rank-$M$ generating twisted
anchor stable under the canonical full order $P$, and put
$B=q_*\mathcal End(A)$. For a punctured twisted object $F_U$, let
$Q$ be the canonical full ordinary pullback of
$j_*\mathcal Hom(A_U,F_U)$. All of these objects are geometric objects
over $k$; none is a varying relative reflexive hull. Use right modules
for the precomposition action. The proved constructions give
\[
 P\subset B\subset P^\vee\subset P(C),\qquad
 Q\subset E=q_*\mathcal Hom(A,F)\subset Q(C),
\]
where
\[
 F=\bigl((q^*Q\otimes_{q^*P}A)_{\rm tf}\bigr)^{\vee\vee}.
\]

\paragraph{What the saturation construction does identify.}
Inside the specified punctured module let $N=Q\cdot B$ denote the
actual image of $Q\otimes_P B$. The exact fixed-anchor Morita functor
identifies this image with the torsion-free image of
$q^*Q\otimes_{q^*P}A$. On a regular two-dimensional root chart its
reflexive hull differs from that image only in codimension two.
At height one the chart is a regular discrete valuation ring, so the
torsion-free module is already free and reflexive. Taking the exact
tame invariant Hom preserves the corresponding image and its
height-one identification. Consequently
\[
                     E=N^{\vee\vee}=(Q\cdot B)^{\vee\vee}
\]
as ordinary sheaves on the smooth coarse surface. Indeed $E$ is
ordinary locally free, $N$ is torsion-free, and $E/N$ has finite
support; these facts characterize the ordinary reflexive hull.
This is an equality for the fixed geometric construction, not a
claim that the image or hull commutes with arbitrary coefficient
base change.

Furthermore $\Gamma(Y,E)=\Gamma(Y,Q)=Q_0$. Both inclusions
$Q\subset E\subset Q(*C)$ respect the punctured identification,
and the canonical $Q$ has the punctured Hartogs equality. This
observation bounds neither $H^1(E)$ nor the length or generators
of the finite-support saturation quotient.

\paragraph{Module generation is a distinct issue.}
The proved ordinary generating quotient for $Q(-Z)$ pulls
back to the refined model. The preceding identity therefore
shows that $E(-Z)$ is generically generated as a $B$-module:
the global generators of $Q(-Z)$ generate $(Q\cdot B)(-Z)$,
and its remaining saturation quotient has finite support.
This uses the actual pulled-back quotient; it does not assert
that an unmodified full bundle remains generically generated
on every newly exceptional component. It gives no small
number of global $B$-generators.

Even for the split algebra $\operatorname{Mat}_d(k)$, the vectors
$v_i=e_i e_1^T$, $1\leq i\leq d$, generate its regular right
module. Every scalar linear combination of these matrices has
image contained in a single line. More generally $b$ such
combinations have combined image of dimension at most $b$.
Thus no choice of $b<d$ scalar combinations generates the right
module. This elementary example distinguishes $B$-generation
from the stronger ordinary evaluation-surjectivity used in the
existing incidence argument. It is not asserted to be a
geometric counterexample with a canonical order.

If $E(-aZ)$ is actually generated as an ordinary bundle, the
known incidence calculation can select $s+1$ sections generating
it as a $B$-module, where $E$ is locally $B^s$ and $s=r/M$.
The presently missing implication is precisely that ordinary
generation, with $a$ independent of $M,t,r$, or another actual
small-source construction replacing it. The rank identity
$M^2s^2=r^2$ alone supplies neither implication.

\paragraph{The precise noncommutative Kahn input.}
Drozd and Gavran (2014), in the author version
arXiv:1301.2999v3, define a weak reduction ideal for a regular
normal order $B$ to be a locally bi-principal ideal $I$ with
exceptional support, such that $I/I^2$ is generically generated
over $B/I$ and $H^1(I/I^2)=0$ (Definition 2.6, p.~7).
Their Corollary 2.9, with its entire proof on p.~8, equivalently
requires generic $B$-generation of $I$ and $H^1(Y,I)=0$.
Their notion of generation is the order-module notion;
Remark 2.2 explicitly distinguishes it from ordinary generation.

Theorem 2.13 and its full written proof on pp.~11--12 give the
pointwise two-level full-module correspondence after this
weak ideal is supplied. Uniqueness uses the additional
generation condition for the dualizing order module. The
proof constructs formal projective modules on two affine
opens and applies proper existence; it does not give a
quantitative all-coefficient polynomial atlas.

Proposition 2.10 obtains an ideal $I=B(-nD)$ by relative
Serre generation and vanishing for a sufficiently positive
exceptional $D$. The integer $n$ is unspecified and may
depend on the order. Consequently this existence theorem
does not prove $H^1(Y,B(-aZ))=0$ for a fixed $a$ in the
degree-$t$, variable-period family. Scalar vanishing for
$\mathcal O_Y(-Z)$ does not imply that order-valued vanishing.
The source's Corollary 2.14 and Theorem 3.3 explicitly omit
their proofs; they remain cited known results and have not
been used to fill this uniformity gap.

\paragraph{Why reduced balanced anchors do not settle the gap.}
The separate C note gives a genuinely balanced residual
anchor with an $O_R(tM^2)$ endomorphism budget on the
\emph{reduced stack divisor}. Removing one coarse exceptional
component instead removes $m_i$ root-normal layers. A layer
of order $j$ includes the corresponding conormal character
and degree shift. Vanishing on the first reduced layer
therefore does not imply vanishing on that full coarse
component, nor a Hartogs equality for the chosen anchor's
endomorphism order. Making the anchor stable under $P$
can change its balanced residual data.

In particular this note proves no uniform $H^1$ tail,
ordinary generation, short order-projective syzygy, finite
object atlas, representation count, or main-gate reduction.
The finite-support saturation identity and the exact
primary-source hypotheses above isolate the remaining
question without replacing it by generic rank or a divisor
sandwich.

\paragraph{Reference and reading scope.}
Drozd, Y. A., \& Gavran, V. S. (2014).
On Cohen--Macaulay modules over non-commutative surface
singularities. \emph{Central European Journal of Mathematics,
12}(5), 675--687.
\url{https://doi.org/10.2478/s11533-013-0374-x}.
The author preprint v3 was opened at
\url{https://imath.kyiv.ua/~drozd/cm-noncommutative.pdf}.
The complete written Sections 1--2, pp.~2--13, were read,
including Theorem 2.1, Lemma 2.5, Propositions 2.7--2.8,
Corollary 2.9, Proposition 2.10, Lemma 2.11,
Proposition 2.12, and Theorem 2.13. Recursively cited
terminal-order existence, maximal-order structure and
proper existence theorems remain known inputs; no new
proof or quantitative scope is attributed to them.
